\section{Part 1.B -- LoRA Fine-Tuning}

\subsection*{Introduction}
This exercise extends the language-model setup by fine-tuning a pre-trained GPT2 model with Low-Rank Adaptation. Instead of updating all parameters, I inserted trainable adapters into the attention projections and kept the base model frozen. The main objective was to achieve a better perplexity than the from-scratch baseline while training only a small fraction of parameters.

\subsection*{Implementation details}
I implemented LoRA manually by modifying the query, key, and value projections inside the GPT2 attention block. The original weights were reused, while the adapter matrices introduced a low-rank update controlled by rank and scaling factor. During training, only the adapter parameters were marked as trainable. This kept the optimization lightweight and made it possible to experiment with different rank/alpha settings without retraining the full model.

\subsection*{Results}
The final version of the report will compare the LoRA configurations against the best from-scratch GPT2 run and highlight the parameter-efficiency trade-off.